\documentclass{article}
\usepackage[T1]{fontenc}
\usepackage{lmodern}
\usepackage{graphicx}
\usepackage{amsmath,amssymb}
\usepackage[a4paper,margin=2cm]{geometry}
\usepackage[absolute,overlay]{textpos}
\setlength{\TPHorizModule}{1mm}
\setlength{\TPVertModule}{1mm}
\usepackage[indonesian]{babel}
\usepackage{hyperref}
\setlength{\emergencystretch}{2em}
% unit: Halaman 21

\begin{document}

% === Halaman 21 (python/native) ===

21

• Hasil perhitungan frekuensi kemunculan huruf, bigram, dan trigram: 

  - huruf I paling sering muncul, 

 - XL adalah bigram yang paling sering muncul, 

  - XLI adalah trigram yang paling sering muncul. 

 Ketiga data terbanyak ini menghasilkan dugaan bahwa

    I  berkoresponden dengan huruf plainteks e, 

    XLI  berkoresponden dengan the, 

    XL berkoresponden dengan th

Pemetaan:

  I → e

  X → t

  L → h

\end{document}
